\usetikzlibrary{decorations.markings}

\begin{tikzpicture}
	
	\defname[named]{strong_enh}{strong enhancers}{35Senhancer}
	\defname[named]{medium_enh}{medium enhancers}{strong_enh}\colorlet{medium_enh}{medium_enh!67!black}
	\defname[named]{weak_enh}{weak enhancers}{strong_enh}\colorlet{weak_enh}{weak_enh!33!black}
	\defname[named]{inactive}{inactive sequences}{gray}
	\defname[named]{repressive}{repressive elements}{Firebrick1}
	
	%%% scheme of category definition %%%
	\coordinate (category definition) at (0, 0);
	
	\tikzset{
		save coordinate/.style = {
			postaction = decorate,
			decoration = {
				markings,
				mark = at position 0.55 with {
					\coordinate (#1);
				}
			}
		},
		mean/.style = {
			densely dashed,
			save coordinate = mean
		},
		mean-1/.style = {
			densely dotted,
			save coordinate = mean-1
		},
		mean+1/.style = {
			densely dotted,
			save coordinate = mean+1
		},
		mean+2/.style = {
			densely dotted,
			save coordinate = mean+2
		},
		mean+3/.style = {
			densely dotted,
			save coordinate = mean+3
		},
		max/.style = {
			draw = none,
			save coordinate = max
		},
		min/.style = {
			draw = none,
			save coordinate = min
		}
	}
	
	\begin{pggfplot}[%
		at = {(category definition)},
		total width = \fourcolumnwidth,
		xlabel = {},
		ylabel = {\enrichment[GC-normalized]},
		ylabel style = {inner xsep = 0pt},
		ylabel add to width = .33\baselineskip,
		enlarge x limits = {rel = 0, abs = .75},
		enlarge y limits = {lower = 0.1},
		xtick = {1, 5},
		xticklabels = {non-ACR\\sequences, all test\\sequences},
		xticklabel style = {align = center},
		col title = light,
		title is csname,
		title color from name,
		axis background/.style = {fill = none},
		tight y
	]{category_def}
	
		\pggf_violin(
			fill = light
		)
	
		\pggf_hline(
			style from column = intercept_id
		)
		
		\pggf_annotate(
			annotation = {
				\node[anchor = east, fill = white, font = \figsmall, xshift = -.5em] at (mean) {mean};
				\draw[{|<[width = 3]}-{>|[width = 3]}, shorten both = -.5\pgflinewidth] (mean) -- (mean-1) node[pos = .5, anchor = west, font = \figsmall] {sd};
				\draw[{|<[width = 3]}-{>|[width = 3]}, shorten both = -.5\pgflinewidth] (mean) -- (mean+1) node[pos = .5, anchor = west, font = \figsmall] {sd};
				\draw[{|<[width = 3]}-{>|[width = 3]}, shorten both = -.5\pgflinewidth] (mean+1) -- (mean+2) node[pos = .5, anchor = west, font = \figsmall] {sd};
				\draw[{|<[width = 3]}-{>|[width = 3]}, shorten both = -.5\pgflinewidth] (mean+2) -- (mean+3) node[pos = .5, anchor = west, font = \figsmall] {sd};
			}
		)
	
	\end{pggfplot}
	
	\begin{pgfonlayer}{background}
		\shade[left color = white, right color = strong_enh!50] (mean+3) -| (max -| pggfplot c1r1.east) node[pos = .75, anchor = west, node font = \figsmall, xshift = .225cm] {strong enhancers} -- (pggfplot c1r1.north east) -| cycle;
		\shade[left color = white, right color = medium_enh!50] (mean+2) -| (mean+3 -| pggfplot c1r1.east) node[pos = .75, anchor = west, node font = \figsmall, xshift = .35cm] {medium enhancers} -| cycle;
		\shade[left color = white, right color = weak_enh!50] (mean+1) -| (mean+2 -| pggfplot c1r1.east) node[pos = .75, anchor = west, node font = \figsmall, xshift = .225cm] {weak enhancers} -| cycle;
		\shade[left color = white, right color = black!50] (mean-1) -| (mean+1 -| pggfplot c1r1.east) node[pos = .75, anchor = west, node font = \figsmall, xshift = .35cm] (inactive) {inactive sequences} -| cycle;
		\shade[left color = white, right color = repressive!50] (mean-1) -| (min -| pggfplot c1r1.east) node[pos = .75, anchor = west, node font = \figsmall, xshift = .225cm] {repressive elements} -- (pggfplot c1r1.south east) -| cycle;
	\end{pgfonlayer}
	
	\draw[decorate, decoration = {brace, amplitude = .15cm, raise = .075cm}] (max -| pggfplot c1r1.east) -- (mean+3 -| pggfplot c1r1.east);
	\draw[decorate, decoration = {brace, amplitude = .15cm, raise = .2cm}] (mean+3 -| pggfplot c1r1.east) -- (mean+2 -| pggfplot c1r1.east);
	\draw[decorate, decoration = {brace, amplitude = .15cm, raise = .075cm}] (mean+2 -| pggfplot c1r1.east) -- (mean+1 -| pggfplot c1r1.east);
	\draw[decorate, decoration = {brace, amplitude = .15cm, raise = .2cm}] (mean+1 -| pggfplot c1r1.east) -- (mean-1 -| pggfplot c1r1.east);
	\draw[decorate, decoration = {brace, amplitude = .15cm, raise = .075cm}] (mean-1 -| pggfplot c1r1.east) -- (min -| pggfplot c1r1.east);


	%%% number of sequences in each category %%%
	\coordinate[xshift = \columnsep] (category numbers) at (category definition -| inactive.east);
	
	\distance{category numbers}{\textwidth, 0}
	
	\begin{pggfplot}[
		at = {(category numbers)},
		total width = \xdistance,
		xlabel = number of test sequences (\times \pgfmathprintnumber{1000}),
		ylabel = {~},
		ylabel style = {inner sep = 0pt},
		ylabel add to width = 34pt,
		yticklabels are csnames,
		title is csname,
		title color from name,
		xmin = 0,
		enlarge x limits = {lower = 0},
		scaled x ticks = manual:{}{\pgfmathparse{##1/1000}},
		y dir = reverse,
		hide legend
	]{category_numbers}
	
		\pggf_bar(
			xbar,
			draw = black,
			style from column = {fill = yticklabels},
			point meta = x,
			nodes near coords,
			nodes near coords align = {horizontal},
			nodes near coords style = {
				text = black,
				opacity = 1,
				font = \figsmall,
				text depth = 0pt,
				/pgf/number format/fixed
			},
			coordinate style/.condition = {x > 150000}{anchor = east}
		)
	
	\end{pggfplot}


	%%% correlation between TF effects and rel. enrichment of TFBSs in sequences grouped by their strength %%%
	\coordinate[yshift = -\columnsep] (correlation) at (category definition |- xlabel.south);
	
	\begin{pggfplot}[%
		at = {(correlation)},
		total width = \textwidth,
		xlabel = {$\log_2$(effect size)},
		ylabel = {$\log_2$(average TFBS counts, normalized to inactive sequences)},
		ylabel add to width = .33\baselineskip,
		xtick = {-1.2, -.9, ..., 1.2},
		title is csname,
		title color from name,
	]{categories_TFs}
	
		\pggf_scatter(color from col name)
	
		\pggf_trendline(/pgfplots/color from row name)
	
		\pggf_stats(stats = {slope, goodness of fit})
	
	\end{pggfplot}


	%%% subfigure labels %%%
	\subfiglabel{category definition}
	\subfiglabel{category numbers}
	\subfiglabel{correlation}

\end{tikzpicture}